% Clase 5 --- Flujo de un campo uniforme a través de una superficie inclinada.
% El vector área A = A n apunta según la normal; Phi = v . A = v A cos(theta).
\begin{center}
\begin{tikzpicture}
  % campo uniforme (flechas horizontales)
  \foreach \y in {-1.2,-0.6,0,0.6,1.2}
     \draw[figvecaux] (-3,\y) -- (-1.4,\y);
  \node[figlblaux] at (-3.15,1.6){$\vec v$};
  % superficie inclinada vista de canto
  \draw[figline,line width=1.7pt] (0,-1.3) -- (1.3,1.3);
  \node[figlbl,figblue] at (1.6,1.4){$A$};
  % normal n
  \draw[figvec] (0.65,0) -- (1.9,-0.63);
  \node[figlbl,figblue,right=1pt] at (1.9,-0.63) {$\hat n$};
  % dirección del campo en la superficie
  \draw[figvecres] (0.65,0) -- (2.25,0);
  \node[figlbl,figred,right=1pt] at (2.25,0) {$\vec v$};
  % ángulo theta entre v y n
  \draw[figaux] (1.35,0) arc (0:-27:0.7);
  \node[figlblaux] at (1.68,-0.24){$\theta$};
\end{tikzpicture}\\[0.4em]
{\small\itshape Flujo a través de una superficie inclinada:
$\Phi=\vec v\cdot\vec A=vA\cos\theta$, con $\vec A=A\,\hat n$ (normal a la cara).}
\end{center}
